%% resumen.tex

% From mitthesis package
% Version: 1.02, 2024/06/19
% Documentation: https://ctan.org/pkg/mitthesis
% Author: Zhenyun Xie

\chapter*{Resumen}
\pdfbookmark[0]{Resumen}{resumen}
La fotónica de silicio ha surgido como una tecnología prometedora para las interconexiones de centros de datos de próxima generación, al permitir comunicaciones ópticas de alto ancho de banda y alta eficiencia energética sobre una plataforma integrada y compacta.
La fotónica integrada programable (Programmable Integrated Photonics, PIP) amplía aún más estas capacidades al permitir circuitos fotónicos flexibles, definidos por software y reconfigurables dentro de una plataforma hardware unificada.
Esta capacidad de reconfiguración hace que la PIP resulte especialmente atractiva para aplicaciones de redes ópticas.
Entre estas aplicaciones, la conmutación óptica ha despertado un interés particular, ya que permite establecer conectividad óptica directa tanto en arquitecturas de red scale-up como scale-out, evitando las conversiones óptico-eléctrico-óptico.

Sin embargo, a medida que continúa creciendo la escala de las redes de los centros de datos, persisten diversos desafíos.
Entre ellos se encuentran el desarrollo de mecanismos de programación eficientes para la PIP que permitan la reconfiguración en tiempo real de distintas funcionalidades de conmutación;
el diseño de arquitecturas hardware de PIP escalables capaces de soportar conmutadores de gran radix;
y la viabilización del despliegue práctico de estos sistemas en entornos de red.

Para abordar estos desafíos, introducimos la teoría de grafos como una abstracción del hardware PIP mediante un grafo ponderado.
Esta abstracción proporciona una representación matemática que permite desarrollar algoritmos de encaminamiento para el control a nivel de circuito, facilita la integración del hardware PIP con marcos de redes definidas por software (SDN) y permite el modelado analítico para el diseño de arquitecturas hardware.

En primer lugar, implementamos algoritmos de encaminamiento sobre un modelo de propósito general de PIP basado en grafos para soportar interconexiones ópticas, multidifusión (multicasting) y conmutación de circuitos, demostrando que diversos circuitos ópticos pueden reconfigurarse automáticamente sobre una única plataforma sin intervención manual.
Para abordar el reto de escalar tejidos de conmutación de gran radix en hardware PIP, la metodología basada en grafos se amplía para modelar y evaluar comparativamente distintas topologías y estrategias de encaminamiento.
El modelo incorpora métricas de la capa física, como la pérdida de inserción, la diafonía (crosstalk) y la relación señal-ruido (SNR), lo que permite evaluar cuantitativamente los compromisos de diseño y proporcionar directrices prácticas para el diseño del hardware.
Finalmente, apoyándonos en este control a nivel de circuito, implementamos un plano de control definido por software mediante interfaces dirigidas por modelos (model-driven) y gNMI para habilitar la comunicación southbound, alineando los recursos ópticos de la Capa 1 con el marco SDN.

Este trabajo demuestra que una metodología basada en grafos permite la programación automatizada y la optimización del diseño de la conmutación fotónica integrada programable, mientras que el plano de control definido por software desarrollado extiende estas capacidades hacia la integración y orquestación a nivel de red.
En conjunto, estas contribuciones establecen una base escalable para integrar la conmutación fotónica por circuitos en sistemas de redes impulsados por inteligencia artificial, favoreciendo interconexiones ópticas flexibles, de alto rendimiento y fácilmente gestionables.

\chapter*{Resum}
\pdfbookmark[0]{Resum}{resum}
La fotònica de silici ha emergit com una tecnologia prometedora per a les interconnexions de centres de dades de nova generació, en permetre una comunicació òptica d'alta amplada de banda i eficient energèticament dins d'una plataforma integrada compacta.
La fotònica integrada programable (PIP) amplia encara més aquesta capacitat en permetre circuits fotònics flexibles, definits per programari i reconfigurables dins d'una plataforma de maquinari unificada.
Aquesta reconfigurabilitat fa que la PIP siga particularment atractiva per a aplicacions de xarxes òptiques.
Entre aquestes aplicacions, la commutació òptica ha despertat un interés especial perquè permet connectivitat òptica directa tant en arquitectures de xarxa scale-up com scale-out, evitant al mateix temps les conversions òptic-elèctric-òptiques.

No obstant això, a mesura que l'escala de les xarxes de centres de dades continua creixent, romanen diversos reptes, entre ells: desenvolupar una programabilitat eficient de la PIP per a admetre la reconfiguració en temps real de diverses funcionalitats de commutació;
dissenyar arquitectures de maquinari PIP escalables capaces de suportar un radix de commutació elevat;
i possibilitar el desplegament pràctic d'aquests sistemes en entorns de xarxa reals.

Per a abordar aquests reptes, introduïm la teoria de grafs per a abstraure el maquinari PIP com un graf ponderat.
Aquesta abstracció proporciona una representació matemàtica que permet el desenvolupament d'algorismes d'encaminament per al control a nivell de circuit, facilita la integració del maquinari PIP amb marcs de xarxa definits per programari, i dona suport al modelatge analític per al disseny d'arquitectures de maquinari.
En primer lloc, implementem algorismes d'encaminament sobre un model PIP de propòsit general basat en grafs per a admetre interconnexions òptiques, multidifusió i commutació de circuits, demostrant que diversos circuits òptics poden reconfigurar-se automàticament en una única plataforma sense intervenció manual.
Per a abordar el repte d'escalar teixits de commutació d'alt radix sobre maquinari PIP, la metodologia basada en grafs s'amplia per a modelar i comparar diferents topologies i estratègies d'encaminament.
El model incorpora mètriques de capa física com la pèrdua d'inserció, la diafonia i la relació senyal-soroll, la qual cosa permet una avaluació quantitativa dels compromisos de disseny i proporciona orientació pràctica per al disseny del maquinari.
Finalment, partint d'aquest control a nivell de circuit, s'implementa un pla de control definit per programari que utilitza interfícies orientades a models i gNMI per a habilitar la comunicació southbound, alineant els recursos òptics de Capa 1 amb el marc SDN.

Aquest treball demostra que una metodologia basada en grafs permet la programació automatitzada i l'optimització del disseny per a la commutació fotònica integrada programable, mentre que el pla de control definit per programari desenvolupat amplia aquestes capacitats a la integració i orquestració a nivell de xarxa.
En conjunt, aquestes contribucions estableixen una base escalable per a integrar la commutació de circuits fotònics en sistemes de xarxa impulsats per IA, donant suport a interconnexions òptiques flexibles, d'alt rendiment i gestionables.
